\BourbakiExerciseGroup

\begin{enumerate}
\def\labelenumi{\arabic{enumi})}
\item
  设 \(\mathbf{K}\) 是特征为 \(p\) 的域， \(n\) 为一个不是 \(p\) 的倍数的整数， \(\mathbf{R}_n(\mathbf{K})\) 为 \(n\) 次单位根域，以及 \(\mathbf{E}\) 为 \(\mathbf{R}_n(\mathbf{K})\) 的一个循环扩张，其次数为 \(n\) （在 \(\mathbf{K}\) 上）；\(\mathbf{E}\) 因此是多项式 \(X^{n}-\alpha\in\mathbb{R}_{n}(\mathbf{K})[\mathbf{X}]\) 的根域（V，第 89 页，定理 4）。要使 E 成为 K 的伽罗瓦扩张，必要且充分的条件是：对每个自同构 \(\tau\in\operatorname{Gal}\left(\mathbb{R}_{n}(\mathbf{K})/\mathbf{K}\right)\) ，存在一个整数 r\textgreater0 和一个元素 \(\gamma\in\mathbb{R}_{n}(\mathbf{K})\) 使得 \(\tau(\alpha)=\gamma^{n}\alpha^{r}\) .
\item
  \begin{enumerate}
  \def\labelenumii{\alph{enumii})}
  \tightlist
  \item
    设 E 是多项式 \(X^{3}-2\in Q[X]\) 的一个分裂扩张。域 E 是 \(\mathrm{R}_{3}(\mathbf{Q})\) 的 3 次循环扩张，但不包含 Q 的 3 次循环扩张。b) 域 \(\mathrm{E}=\mathrm{R}_{9}(\mathbf{Q})\) 是 \(\mathrm{R}_{3}(\mathbf{Q})\) 的 3 次循环扩张，并且包含 Q 的 3 次循环扩张。
  \end{enumerate}
\end{enumerate}

\begin{enumerate}
\def\labelenumi{\alph{enumi})}
\setcounter{enumi}{2}
\tightlist
\item
  在 V，第 154 页，习题 7 的假设和记号下，证明 \(\mathrm{K}(\alpha,\beta)\) 在 K 上为伽罗瓦扩张的充要条件是 d 是 \(\mathrm{K}(\alpha)\) 中的平方；\(\mathrm{K}(\beta)\) 在 K 上为循环扩张的充要条件是 d 是 K 中的平方。
\end{enumerate}

* 3) 证明：对每个整数 \(n > 1\) 都存在一个次数为 \(n\) 的、关于域 \(\mathbf{Q}\) 的循环扩张。（化归到 \(n\) 是一个素数 \(q\) 的情形，并利用乘法群 \((\mathbf{Z} / q^m\mathbf{Z})^*\) 的结构（VII，第 13 页）。*

\begin{enumerate}
\def\labelenumi{\arabic{enumi})}
\setcounter{enumi}{3}
\tightlist
\item
  设 \(\mathbf{K}\) 是特征为 \(p, q\) 一个素数 \(\neq p\) ，使得 \(\mathbf{K}\) 包含 \(q\) 次单位根，并设 \(\zeta\) 是一个本原 \(q\) 次单位根。
\end{enumerate}

\begin{enumerate}
\def\labelenumi{\alph{enumi})}
\item
  设 E 是 K 的循环扩张，次数为 \(q^e\) ，且 \(\sigma\) 是 E 的一个生成 Gal(E/K) 的 K-自同构。设 F 是 K 上次数为 \(q^{e-1}\) 的域，使得 \(K \subset F \subset E\) ；存在 \(\theta \in E\) ，为某个不可约多项式 \(X^q - \alpha\) （在 F{[}X{]} 中）的一个根，使得 \(E = F(\theta)\) 且 \(\theta^{\sigma^m} = \zeta\theta\) （其中 \(m = q^{e-1}\) ）。证明同样有 \(E = K(\theta)\) 且 \(\theta^\sigma = \beta\theta\) ，其中 \(\beta \in F\) 使得 \(\alpha^{\sigma-1} = \beta^q\) 且 \(N_{F/K}(\beta) = \zeta\) （观察到 \(\theta^\sigma = \beta\theta^k\) ，其中 \(0 < k < q\) 且 \(\beta \in F\) ，由习题 1；通过计算 \(\theta^{\sigma^{mq}}\) ，证明 \(k^{mq} - 1 \equiv 0 \pmod{q}\) ）并推断出 \(k = 1\) ).
\item
  反之，设 F 是 K 的一个次数为 \(q^{e-1}\) 的循环扩张，其中 \(e \geqslant 2\) ，并且令 \(\sigma\) 是 F 的一个生成 Gal(F/K) 的 K-自同构。若存在一个元素 \(\beta \in F\) 使得 \(N_{F/K}(\beta) = \zeta\) ，并且 \(\alpha \in F\) 使得 \(\alpha^{\sigma-1} = \beta^q\) ，证明对每一个 \(c \in K^*\) ，多项式 \(X^q - c\alpha\) 在 F{[}X{]} 中不可约。若 \(\theta\) 是该多项式在 F 的一个代数闭包 \(\Omega\) 中的一个根，证明存在一个 K-同态 \(\bar{\sigma}\) ，它把 \(E = F(\theta)\) 映到 \(\Omega\) 内，扩张 \(\sigma\) ，并且满足 \(\theta^{\sigma} = \beta\theta\) ；由此推出 E 是 K 的次数为 \(q^e\) 的循环扩张，元素 \(\bar{\sigma}\) 生成 Gal(E/K)，并且 \(E = K(\theta)\) 。最后证明 K 的每个包含 F 的次数为 \(q^e\) 的循环扩张都是某个多项式 \(X^q - c\alpha\) （在 F{[}X{]} 中的，取适当的 \(c \in K^*\) ）的根域（利用 V 中第 85 页的定理 3）。
\item
  令 K 为有理数域 Q。多项式 \(X^{2} + 1\) 在 Q{[}X{]} 中不可约；如果 i 是它的一个根，则 \(F = \mathbf{Q}(i)\) 是 Q 的 2 次循环扩张，但不存在包含 F 的 Q 的 4 次循环扩张。
\end{enumerate}

\begin{enumerate}
\def\labelenumi{\arabic{enumi})}
\setcounter{enumi}{4}
\item
  设 K 是特征为 p 的域，n 是不为 p 的倍数的整数。对于 K{[}X{]} 中形如 \(X^{n} - a\) 的一个多项式，要使其不可约，必要且充分条件是：对于 n 的每个素因子 q，a 不等于 K 中某个元素的 q 次幂，并且当 \(n \equiv 0 \pmod{4}\) 时，a 不应具有形式 \(-4c^{4}\) （其中 \(c \in K\) ）。（为证明条件的充分性，化归到 \(n = q^{e}\) （q 为素数）的情形，并利用 V 章第 153 页习题 4。对 e 进行归纳论证，借助 V 章第 153 页习题 4 确定 \(X^{q^{e}} - a\) 的每个不可约因子的常数项的形式。）
\item
  设 \(\mathbf{K}\) 是特征为 \(p, n\) 一个不是 \(p\) 的倍数的整数，且 \(\mathbf{N}\) 为多项式 \(\mathbf{X}^n - a\) （ \(\mathbf{K}[\mathbf{X}]\) 的）的一个分裂域.
\end{enumerate}

\begin{enumerate}
\def\labelenumi{\alph{enumi})}
\tightlist
\item
  证明群 \(\operatorname{Gal}(\mathbf{N} / \mathbf{K})\) 同构于群 \(\Gamma\) （由矩阵 \(\left( \begin{array}{cc}x & y\\ 0 & 1 \end{array} \right)\) ，其中 \(x\in (\mathbf{Z} / n\mathbf{Z})^*\) 且 \(y\in \mathbf{Z} / n\mathbf{Z}\) ，所组成的）的一个子群。若 \(n\) 是素数，群
\end{enumerate}

Gal(N/K) 仅当 \(\mathrm{R}_{n}(\mathrm{K})=\mathrm{K}\) 或者 a 是 K 中某个元素的 n 次幂时才是交换的。b) 如果 K=Q 且 n 是素数，证明若 a 不是某个有理数的 n 次幂，则 Gal(N/K) 同构于整个 \(\Gamma\) .

\begin{enumerate}
\def\labelenumi{\alph{enumi})}
\setcounter{enumi}{2}
\tightlist
\item
  如果 \(\mathbf{K} = \mathbf{Q}\) ，确定 \(\operatorname{Gal}(\mathbf{N} / \mathbf{K})\) （当 \(a\) 是一个无平方因子整数时）。
\end{enumerate}

\begin{enumerate}
\def\labelenumi{\arabic{enumi})}
\setcounter{enumi}{6}
\tightlist
\item
  域 K 的一个扩张 E 称为库默尔扩张，如果存在整数 n \textgreater{} 0 使得 K 包含一个本原 n 次单位根，且 E 是伽罗瓦群被 n 零化的阿贝尔扩张。
\end{enumerate}

\begin{enumerate}
\def\labelenumi{\alph{enumi})}
\item
  设 K 为一个域，N 为 K 的 m 次伽罗瓦扩张，其伽罗瓦群为 \(\Gamma\) 。要使 N 成为 K 的库默尔扩张，必要且充分条件是存在 N 在 K 上的一个基 \((\theta_{\sigma})_{\sigma\in\Gamma}\) ，使得元素 \(\gamma_{\sigma\tau}=\theta_{\sigma}^{1-\tau}\) 对每一对 \((\sigma,\tau)\) 都属于 K。（对于条件的必要性，可归结为 N 是次数为素数幂的循环扩张的情形。对于充分性，注意若 \(\tau\in\Gamma\) 的阶为 n，则存在 \(\sigma\in\Gamma\) 使得 \(\gamma_{\sigma\tau}\) 是本原 n 次单位根。）
\item
  如果 \((\theta_{\sigma})_{\sigma \in \Gamma}\) 是满足条件 \(a)\) 的基，证明对于一个元素 \(x = \sum_{\sigma} a_{\sigma} \theta_{\sigma}\) （ \(N\) 的）
\end{enumerate}

使得 \(x\) 的共轭元构成 \(\mathbf{N}\) 在 \(\mathbf{K}\) 上的一个正规基，其充要条件是 \(a_{\sigma} \neq 0\) 对所有 \(\sigma \in \Gamma\) 成立.

* 8) 设 N 是域 K 的一个库默尔扩张。

\begin{enumerate}
\def\labelenumi{\alph{enumi})}
\item
  设 E 是 N 的一个子扩张，使得 \([E:K]=q\) 是一个素数。证明存在将 N 分解为 K 的循环扩张的分解 \(N=N_{1}\otimes N_{2}\otimes\cdots\otimes N_{r}\) ，使得 \(E\subset N_{j}\) 对某个索引 j 成立。
\item
  由 a) 可推出，如果 \(x \in N\) 使得 x 在 K 上的共轭元素构成 N 在 K 上的正规基，则对于每个满足 \(K \subset E \subset N\) 的域 E，x 在 E 上的共轭元构成 N 在 E 上的一个正规基（归结为 \([E : K]\) 是素数的情形，并利用习题 7，b）。*
\end{enumerate}

\begin{enumerate}
\def\labelenumi{\arabic{enumi})}
\setcounter{enumi}{8}
\tightlist
\item
  设 \(\mathbf{K}\) 是特征为 \(p > 0\) 的域.
\end{enumerate}

\begin{enumerate}
\def\labelenumi{\alph{enumi})}
\item
  设 E 是 K 的循环扩张，次数为 \(p^e\) ，并设 \(\sigma\) 是生成群 \(\operatorname{Gal}(\mathrm{E} / \mathrm{K})\) 的 K-自同构。设 F 是 K 上次数为 \(p^{e - 1}\) 的域，使得 \(\mathbf{K} \subset \mathbf{F} \subset \mathbf{E}\) ；如果 \(m = p^{e - 1}\) ，证明存在 \(\theta \in \mathbf{E}\) ，它是不可约多项式 \(\mathbf{X}^p - \mathbf{X} - \alpha\) （ \(\mathrm{F}[\mathbf{X}]\) 中的）的一个根，使得 \(\mathrm{E} = \mathrm{F}(\theta)\) 且 \(\sigma^m (\theta) = \theta + 1\) 。证明我们还有 \(\mathrm{E} = \mathrm{K}(\theta)\) 且 \(\sigma (\theta) = \theta + \beta\) ，其中 \(\beta \in \mathrm{F}\) 使得 \(\sigma (\alpha) - \alpha = \beta^p - \beta\) 且 \(\mathrm{Tr}_{\mathrm{F} / \mathrm{K}}(\beta) = 1\) .
\item
  反之，设 F 是 K 的次数为 \(p^{e-1}\) 的循环扩张（其中 e \textgreater{} 1), \(\sigma\) 是 F 的一个生成 Gal(F/K) 的 K-自同构。证明存在 F 的两个元素 \(\alpha\) , \(\beta\) ，使得 \(\operatorname{Tr}_{\mathrm{F}/\mathrm{K}}(\beta) = 1\) 且 \(\sigma(\alpha) - \alpha = \beta^{p} - \beta\) （利用 V 第 85 页定理 3 及 V 第 49 页命题 1 的推论）。由此推出对所有 \(c \in K\) ，多项式 \(X^{p} - X - \alpha - c\) 在 F{[}X{]} 中不可约；如果 \(\theta\) 是该多项式在 F 的一个代数闭包 \(\Omega\) 中的一个根，证明存在一个 K-同态 \(\bar{\sigma}\) ，它把 E = F( \(\theta\) ) 映到 \(\Omega\) 内，延拓 \(\sigma\) 并且满足 \(\bar{\sigma}(\theta) = \theta + \beta\) ；由此得出 E 是 K 的一个次数为 \(p^{e}\) 的循环扩张， \(\bar{\sigma}\) 生成 Gal(E/K) 且 E = K( \(\theta\) )。最后证明，K 的每个包含 F 的 \(p^{e}\) 次循环扩张都是多项式 \(X^{p} - X - \alpha - c\) （F{[}X{]} 中的，取适当的 \(c \in K\) ）的分裂域.
\end{enumerate}

\begin{enumerate}
\def\labelenumi{\arabic{enumi})}
\setcounter{enumi}{9}
\tightlist
\item
  \begin{enumerate}
  \def\labelenumii{\alph{enumii})}
  \tightlist
  \item
    设 \(\mathbf{K}\) 是一个域，其代数闭包是素数次数 \(q\) 的 \(\mathbf{K}\) 的扩张。这样 \(\mathbf{K}\) 是完全域，并且 \(q\) 必须不同于 \(\mathbf{K}\) 的特征（使用习题 9）。
  \end{enumerate}
\end{enumerate}

\begin{enumerate}
\def\labelenumi{\alph{enumi})}
\setcounter{enumi}{1}
\tightlist
\item
  证明 K 包含 q 次单位根，并且 \(\Omega\) 是不可约多项式 \(X^{q}-a\) （K{[}X{]} 中的）的分裂域；推断出 q=2（在相反情形下，由习题 5 推断该多项式 \(X^{q^{2}}-a\) 将是不可约的）。进一步证明 -a 在 K 中是平方（见上述引文），-1 在 K 中不是平方，并由此推断出 \(\Omega=\mathrm{K}(i)\) ，其中 \(i^{2}=-1\) .
\end{enumerate}

\begin{enumerate}
\def\labelenumi{\arabic{enumi})}
\setcounter{enumi}{10}
\item
  设 \(\mathbf{K}\) 是一个域，其代数闭包 \(\Omega\) 是一个有限次 \(> 1\) 的 \(\mathbf{K}\) 的扩张。证明 \(\Omega = \mathbf{K}(i)\) ，其中 \(i^2 = -1\) 。（如果我们有 \(\Omega \neq \mathbf{K}(i)\) ，借助伽罗瓦理论和西罗定理证明，那么就会存在一个域 \(E\) 使得 \(\mathbf{K}(i) \subset E \subset \Omega\) 且 \(\Omega\) 是 \(E\) 的素数次扩张；现在应用习题 10。）
\item
  设 K 是一个特征为 q 的域， \(\Omega\) 为 K 的一个代数闭扩张，x 为 \(\Omega\) 的一个元素，使得 \(x \notin K\) ，且 M 为 \(\Omega\) 的不包含 x 的子扩张所成集合的一个极大元（V，第 156 页，习题 9）；进一步假设 M 是完全域，使得 \(\mathbf{M}(x)\) 是素数次数 p 的 M 的伽罗瓦扩张（同上）。
\end{enumerate}

\begin{enumerate}
\def\labelenumi{\alph{enumi})}
\item
  证明：要么 \(p = 2\) 且 \(\Omega = M(i)\) ，要么对于每个整数 \(r \geqslant 1\) ，存在唯一一个次数为 \(p^r\) 的、关于 \(M\) 的扩张；该扩张 \(M_r\) 在 \(M\) 上是循环的， \(\Omega\) 是诸 \(M_r\) 的并，且诸 \(M_r\) 是 \(M\) 的仅有的有限次扩张。（注意，如果一个 \(p\) -群 \(\Gamma\) 的阶为 \(p^r\) ，且它只有一个阶为 \(p^{r-1}\) 的子群，则它必然是循环的；可对 \(r\) 作归纳论证，并利用 I，第 133 页，习题 10，以及 I，第 77 页，命题 11 的推论）。
\item
  假设 \(\Omega \neq \mathbf{M}(i)\) 且 \(p \neq q\) 。证明 \(\mathbf{M}\) 包含 \(p\) 次单位根；于是我们有 \(\mathbf{M}(x) = \mathbf{M}_1 = \mathbf{M}(\alpha)\) ，其中 \(\alpha^p \in \mathbf{M}\) 。证明若 \(\alpha\) 是一个 \(p\) 次幂（ \(\mathbf{M}_1\) 中某个元素的），则 \(\mathbf{M}_1 = \mathbf{M}(\zeta)\) ，其中 \(\zeta\) 是一个 \(p^2\) 次单位根（若 \(\alpha = \beta^p\) ，其中 \(\beta \in \mathbf{M}_1\) ，考虑 \(\beta\) 在 \(\mathbf{M}[\mathbf{X}]\) 中的极小多项式）。
\item
  在 \(b\) ) 的假设下，证明若 \(\alpha\) 不是一个 \(p\) 次幂（ \(\mathbf{M}_1\) 中某个元素的），则 \(\mathbf{M}_r\) 是多项式 \(\mathbf{X}^{p^{r-1}} - \alpha\) （ \(\mathbf{M}_1[\mathbf{X}]\) 中的）的分裂域（利用 \(b\) )).
\end{enumerate}

\(d)\) 在 \(b\) ) 的假设下，证明若 \(\alpha\) 是一个 \(p\) 次幂（ \(\mathbf{M}_1\) 中某个元素的），则存在 \(\gamma \in \mathbf{M}_1\) 使得 \(\gamma \notin \mathbf{M}\) 且 \(\mathbf{M}_2\) 是 \(\mathbf{X}^p - \gamma\) 的分裂域；证明 \(\gamma\) 不是一个 \(p\) 次幂（ \(\mathbf{M}_2\) 中某个元素的），并由此推出 \(\mathbf{M}_r\) 是多项式 \(\mathbf{X}^{p^{r-1}} - \gamma\) （ \(\mathbf{M}_1[\mathbf{X}]\) 中的）的分裂域。

\begin{enumerate}
\def\labelenumi{\arabic{enumi})}
\setcounter{enumi}{12}
\tightlist
\item
  如果域 K 的每一个在 K 上的有限次可分代数扩张都是循环扩张，则称 K 为 Moriya 域。* 每个有限域都是 Moriya 域（V，第 95 页，命题 4）。* 习题 12 中的域 M 是 Moriya 域。
\end{enumerate}

\begin{enumerate}
\def\labelenumi{\alph{enumi})}
\item
  要使 K 成为 Moriya 域，当且仅当对每个整数 n \textgreater{} 0，在 K 上至多存在一个 n 次可分扩张。（使用如下事实：若 G 是有限群，且对每个整除 G 的阶的整数 \(m \geqslant 1\) ，G 中至多存在一个 m 阶子群，则 G 是循环群；这可由 p-群的同一性质（习题 12）以及 I，第 80 页，定理 4 推出。）
\item
  若 K 是 Moriya 域，证明 K 的每个代数扩张 E 都是 Moriya 域，并且若 F 是 E 上的一个 m 次扩张，则存在 K 上的一个 m 次扩张。
\end{enumerate}

\begin{enumerate}
\def\labelenumi{\arabic{enumi})}
\setcounter{enumi}{13}
\tightlist
\item
  \begin{enumerate}
  \def\labelenumii{\alph{enumii})}
  \tightlist
  \item
    设 K 是一个域，使得存在 K 的任意大有限次的代数扩张，且所有这些扩张的次数都是同一个素数 p 的幂。证明在这些条件下，对每个幂 \(p^{h}\) ( \(h \geqslant 0\) ）都存在 K 的一个次数为 \(p^{h}\) 的扩张（利用 I，第 77 页，定理 1）。
  \end{enumerate}
\end{enumerate}

\begin{enumerate}
\def\labelenumi{\alph{enumi})}
\setcounter{enumi}{1}
\tightlist
\item
  设 p 为一个素数，K 为一个域，使得存在 K 的一个代数扩张，其次数在 \(p \neq 2\) 时被 p 整除，在 p = 2 时被 4 整除。证明存在 K 的一个代数扩张 E，使得 E 的代数扩张的次数的集合恰为 p 的幂 \(p^{h}\) 的集合。（若 K 是完全域或特征 \(\neq p\) ，令 \(K_{1}\) 为 K 的完全闭包；考虑 \(K_{1}\) 的一个极大的代数扩张，在所有其元素在 \(K_{1}\) 上的次数不被 p 整除的那些代数扩张之中选取；利用 a) 和习题 11。）
\end{enumerate}

\begin{enumerate}
\def\labelenumi{\arabic{enumi})}
\setcounter{enumi}{14}
\item
  设 K 为 Moriya 域（习题 13），P 为整除 K 的有限次代数扩张的次数的素数集合。证明：K 的有限次代数扩张的次数的集合，要么等于集合 \(N_{P}\) （即其所有素因子均在 P 中的整数所构成的集合），要么等于 \(N_{P}\) 的由不能被 4 整除的整数组成的子集。
\item
  对于每个素数 \(p\) ，存在一个特征为 \(p\) 的 Moriya 域，使得该域的有限次代数扩张的次数的集合是集合 \(\mathbf{N}^*\) （即所有 \(>0^1\) 的整数所组成的集合）。证明对于素数的每一个集合 \(P\) ，存在一个任意特征的 Moriya 域 \(K\) ，使得 \(K\) 的有限次代数扩张的次数的集合是集合 \(N_P\) （如习题 14, b) 中那样论证）。
\end{enumerate}

* 17) 设 \(\mathbf{P}_n\) 为首一多项式 \(\mathbf{X}^n + a_1\mathbf{X}^{n-1} + \cdots + a_n\) （ \(\mathbf{Z}[\mathbf{X}]\) 中的）所成的集合，这些多项式在 \(\mathbf{C}\) 中的所有根都具有绝对值 \(\leqslant 1\) .

\begin{enumerate}
\def\labelenumi{\alph{enumi})}
\item
  证明集合 \(P_{n}\) 是有限的。由此推出，对于每个多项式 \(F \in P_{n}\) ，如果 \(F_{h}\) 表示以F的根的h次幂为根的多项式，则存在两个不同的整数h、k，使得 \(F_{h} = F_{k}\) .
\item
  由 a) 推出，多项式 \(F \in P_{n}\) 的所有根要么为零，要么是单位根（“克罗内克定理”）。*
\end{enumerate}

\begin{enumerate}
\def\labelenumi{\arabic{enumi})}
\setcounter{enumi}{17}
\tightlist
\item
  设 \(p \in \mathbb{N}\) 是一个素数， \(C_p\) 为级为 \(p\) 的、关于 \(\mathbf{Q}\) 的分圆扩张， \(\zeta \in C_p\) 为一个本原 \(p\) 次单位根， \(n\) 为一个整数。记
\end{enumerate}

\[
\mathrm{X} _ {n} (\zeta) = \sum_ {\alpha \in \mathbf {Z} / p \mathbf {Z}} \zeta^ {\alpha^ {n}}.
\]

\begin{enumerate}
\def\labelenumi{\alph{enumi})}
\item
  计算 \(d = [\mathbf{Q}(\mathbf{X}_n(\zeta)) : \mathbf{Q}]\) 。确定扩张 \(\mathbf{Q}(\mathbf{X}_n(\zeta))\) （ \(\mathbf{Q}\) 上的）的伽罗瓦群，以及 \(\mathbf{X}_n(\zeta)\) 的共轭元。
\item
  设 \(T^{d} + a_{1}T^{d-1} + \cdots + a_{d}\) 为 \(\mathbf{X}_{n}(\zeta)\) 在 Q 上的极小多项式。证明 \(a_{1} = 0\) .
\item
  假设 \(n > 1\) 且 \(p \equiv 1 \pmod{n}\) 。证明 \(a_2 = -\frac{n(n - 1)p}{2}\) （当 \(\frac{p - 1}{n}\) 为偶数时），而 \(a_2 = \frac{np}{2}\) （当 \(\frac{p - 1}{n}\) 为奇数时）。
\item
  假设 \(p \equiv 1 \pmod{n}\) 。设 \(\beta \in \mathbf{F}_p^*\) 是这样一个元素，其在 \(\mathbf{F}_p^* / \mathbf{F}_p^{*n}\) 中的剩余类是一个生成元。设 \(\mu\) 为方程在 \(\mathbf{F}_p^n\) 中的解的个数，该方程为 \(\sum_{i=1}^{n} \beta^{i-1} \mathbf{X}_i^n = 0\) 。证明 \(a_n = \frac{(-1)^n p(\mu - p^{n-1})}{p-1}\) .
\item
  计算 \(a_3\) ，其中 \(n = 3\) 且 \(p = 7, 13, 19\) .
\end{enumerate}

*1 域 \(\mathbf{Q}_p\) （即 \(p\) -进数域）在其代数闭包中的最大完全分歧扩张是一个具有相同性质的特征为 0 的域。*

* 19) a) 假设素数定理成立；如果 \(\pi(x)\) 是满足 \(p \leqslant x\) 的素数的个数，则 \(\pi(x) \sim x / \log x\) （当 \(x\) 趋于 \(+\infty\) ）。由此推出，对于每个奇数 \(m\) ，存在一个递增的素数序列 \(p_1 < p_2 < \cdots < p_m\) 使得 \(p_m < p_1 + p_2\) （考虑介于 \(x\) 与 \(3x/2\) 之间的素数）。

\begin{enumerate}
\def\labelenumi{\alph{enumi})}
\setcounter{enumi}{1}
\tightlist
\item
  利用由 a) 确定的有限素数序列 \(p_{j}\) ( \(1 \leq j \leq m\) ） ，令 \(n = p_{1}p_{2} \ldots p_{m}\) 。证明在分圆多项式 \(\Phi_{n}(X)\) 中，次数为 \(< p_{m} + 1\) 的项与多项式
\end{enumerate}

\[
(1 - \mathbf {X}) ^ {- 1} \prod_ {j = 1} ^ {m} (1 - \mathbf {X} ^ {p _ {j}}).
\]

的相应项相同。由此推出 \(\mathbf{X}^{p_m}\) 在 \(\Phi_n(\mathbf{X})\) 中的系数是 \(1 - m\) （“I. Schur 定理”）。*

* 20) 设 \(p\) 是一个素数，且 \(G\) 为一个交换 \(p\) -挠群（VII，第 10 页）。\\
a) 设 \(G_i \subset G\) 为满足以下条件的元素 \(x \in G\) 的集合：对每个 \(l\) ，方程 \(p^{l} y = x\) 有一个解。证明 \(G_i\) 是 \(G\) 的一个直因子。

\begin{enumerate}
\def\labelenumi{\alph{enumi})}
\setcounter{enumi}{1}
\item
  假设方程 px = 0 在 G 中只有有限多个解。证明：对每个整数 l，核 \(_{l}G\) （即 G 中乘以 \(p^{l}\) 的核）是一个有限群，并且秩 \(r(l)\) （即 \(F_{p}\) -向量空间 \(_{l}G/_{(l-1)}G\) 的秩）是 l 的递减函数，因此当 \(l > l_{0}\) （ \(l_{0}\) 足够大）时为常数。（利用初等因子定理（VII，第 24 页）。证明 \(G/G_{i}\) 是一个被 \(p^{l_{0}}\) 零化的有限群。）
\item
  进一步假设 \(\mathbf{G} = \mathbf{G}_i\) 。证明映射 \(l \mapsto r(l)\) 是常数。记 \(\mathbf{H} = (\mathbf{Z}[1/p]/\mathbf{Z})^{(1)}\) ，证明对每个 \(l\) 存在同构 \(_l\mathbf{G} \to {}_l\mathbf{H}\) ，并且每个同构 \(_l\mathbf{G} \to {}_l\mathbf{H}\) 扩展为同构 \((l+1)\mathbf{G} \to {}_{(l+1)}\mathbf{H}\) 。由此推出 \(\mathbf{G}\) 同构于 \(\mathbf{H}\) .
\item
  在 b) 的假设下，证明存在整数 n 和一个映射 \(i \to n_i\) （从 \(N - \{0\}\) 到 N 中，具有有限支撑），使得 G 同构于 \((\mathbf{Z}[1/p]/\mathbf{Z})^n \oplus \bigoplus_{i} (\mathbf{Z}/p^i \mathbf{Z})^{n_i}\) 。证明整数 n 和映射 \(i \mapsto n_i\) 是唯一地
\end{enumerate}

由映射 \(l \mapsto \text{Card}(l\mathbf{G})\) 所决定的。*

* 21) 设 G 和 H 是两个交换挠群，使得对每个整数 \(n\) ，方程 \(nx = 0\) 在 G 和 H 中具有相同的有限个解。证明 G 与 H 同构。（将 G 和 H 分解为 \(p\) -挠群的直和并应用习题 20。）令 \(G = \bigoplus_{i \in \mathbf{N}} \mathbf{Z} / p^i\mathbf{Z}\) 且 \(H = \bigoplus_{i \in \mathbf{N}} \mathbf{Z} / p^{2i}\mathbf{Z}\) 。证明 G 与 H 不同构

并且对每个整数 \(n\) ，方程 \(nx = 0\) 的解集的基数在 \(G\) 和 \(H\) 中是相同的。*

\begin{enumerate}
\def\labelenumi{\arabic{enumi})}
\setcounter{enumi}{21}
\tightlist
\item
  \begin{enumerate}
  \def\labelenumii{\alph{enumii})}
  \tightlist
  \item
    设 k 是一个域，f 是 \(k[X]\) 中的一个可分的首一多项式，D 是 f 的一个分裂扩张， \(f, x_{1}, \ldots, x_{n}\) 是 f 在 D 中的根，且 G 是 D 在 k 上的伽罗瓦群。假设 G 是循环群。证明：要使 G 的每个元素都诱导 \(x_{i}\) 的一个偶置换，其必要且充分条件是 f 的不可约因子的个数与 f 的次数模 2 同余（考虑 G 的一个生成元 \(\sigma\) 以及 \(x_{i}\) 的置换（由 \(\sigma\) 所诱导的）的轮换分解）。
  \end{enumerate}
\end{enumerate}

\begin{enumerate}
\def\labelenumi{\alph{enumi})}
\setcounter{enumi}{1}
\tightlist
\item
  进一步假设 k 的特征为 \(\neq 2\) 。证明上述两个条件等价于 f 的判别式是 k 中的平方这一事实（“Stickelberger 定理”）。
\end{enumerate}

\begin{enumerate}
\def\labelenumi{\arabic{enumi})}
\setcounter{enumi}{22}
\tightlist
\item
  设 \(n\) 和 \(p\) 是两个奇且互异的素数。
\end{enumerate}

\begin{enumerate}
\def\labelenumi{\alph{enumi})}
\item
  设 \(Q \in F_{p}[X]\) 为分圆多项式 \(\Phi_{n}\) 的像。证明：Q 的不可约因子的个数 i 为偶数当且仅当 p 是模 n 的平方。（设 K 为 Q 的分裂域。证明 \(\operatorname{Gal}(K/F_{p})\) 可等同于 \((\mathbf{F}_{n})^{*}\) 的由 p 的像生成的子群。研究指标 \(((\mathbf{F}_{n})^{*}: \operatorname{Gal}(K/F_{p}))\) 的奇偶性。）
\item
  证明 \(\mathbf{X}^n - 1 \in \mathbf{F}_p[\mathbf{X}]\) 的判别式是
\end{enumerate}

\[
(- 1) ^ {\frac {n (n - 1)}{2} + n - 1} n ^ {n}.
\]

\begin{enumerate}
\def\labelenumi{\alph{enumi})}
\setcounter{enumi}{2}
\tightlist
\item
  证明： \(i\) 是偶数当且仅当 \((-1)^{\frac{n(n - 1)}{2}} n\) 是模 \(p\) 的平方（把习题 22 应用于多项式 \(\mathbf{X}^n - 1\) ）。
\end{enumerate}

\[
\left(\frac {n}{p}\right) = 1
\]

\[
\left(\frac {n}{p}\right) = - 1
\]

\(\left(\frac{n}{p}\right)\left(\frac{p}{n}\right) = (-1)^{\frac{(n - 1)(p - 1)}{4}}\) （“二次互反律”）。
% semantic-fragments: legacy-input-seam
\relax{}
